\begin{problem}[第34届物理决赛][双星系统]{81}
    双星系统是一类重要的天文观测对象。假设某两星体均可视为质点，其质量分别为 $M$ 和 $m$，一起围绕它们的质心做圆周运动，构成一双星系统，观测到该系统的转动周期为 $T_{0}$。在某一时刻，$M$ 星突然发生爆炸而失去质量 $\Delta M$。假设爆炸是瞬时的、相对于 $M$ 星是各向同性的，因而爆炸后 $M$ 星的残余体 $M'$ ($M' = M - \Delta M$) 星的瞬间速度与爆炸前瞬间 $M$ 星的速度相同，且爆炸过程和抛射物质 $\Delta M$ 都对 $m$ 星没有影响。已知引力常量为 $G$，不考虑相对论效应。
    \begin{subproblem}
        \begin{subsubproblem}
            求爆炸前 $M$ 星和 $m$ 星之间的距离 $r_{0}$;
        \end{subsubproblem}
        \begin{subsubproblem}
            若爆炸后 $M'$ 星和 $m$ 星仍然做周期运动，求该运动的周期 $T_{1}$;
        \end{subsubproblem}
        \begin{subsubproblem}
            若爆炸后 $M'$ 星和 $m$ 星最终能永远分开，求 $M$、$m$ 和 $\Delta M$ 三者应满足的条件。
        \end{subsubproblem}
    \end{subproblem}
\end{problem}